\section{从事务到关系：memory、warm start 与长期上下文}

讲者指出当前大量 Agent 仍是 transactional：每次会话像第一次见面，缺少跨会话记忆，导致用户重复提供身份信息和问题背景。为此，Sierra 把“warm start”作为关键能力：下一次交互从“第二垒/第三垒”起步，而不是回到“你是谁”。

memory 不是“把所有对话都塞进上下文窗口”。在企业服务中，memory 更接近“结构化客户状态层”：包含身份验证历史、设备/订单状态、近期交互结果、风险标记、可复用偏好等。模型只消费与当前任务相关的“可证明有用”片段。

\begin{importantbox}{warm start 的业务价值}
warm start 直接影响三项指标：平均处理时长（AHT）、首次解决率（FCR）、升级率（escalation）。当 Agent 能在会话开始时拿到有效状态，后续推理链会显著缩短，错误分支也会减少。
\end{importantbox}

\begin{warningbox}{memory 设计的边界}
记忆越多不一定越好。过度保留无关上下文会引入隐私风险和检索噪声，导致模型“看得太多、答得更差”。必须结合 retention policy、数据分级、用途限定进行治理。
\end{warningbox}

\subsection{本章小结}

memory 的目标不是“记住一切”，而是“在合法边界内记住能提升解决率的最小必要信息”。warm start 是 customer-facing agent 从短期事务走向长期关系的基础设施。

